\section{Introduction}
\label{sec:intro}

\noindent
Recent years have seen rapid growth in applications powered
by deep learning models like large language models (LLMs)~\cite{aws-rekognition,chatgpt,copilot,sora,stability-ai}.
Due to the huge computation requirements,
these models are typically served in model-serving-as-a-service systems ({\maas})~\cite{serverless-together-ai,serverless-ray,deepinfra,DBLP:conf/asplos/YangZLZLZCL22,ali-serverless-llm,aws-serverless-llm,DBLP:journals/corr/abs-2401-14351},
which manage a cluster of accelerators (e.g., GPUs) 
and provision an appropriate number of serving \emph{instances} containing GPUs
to each model deployed.
%Specifically, a serving instance is the set of
%GPU(s) used that hold a full copy of a model's parameters.


An {\maas} system has two design objectives:
\emph{maximizing goodput}---the number of requests
that meet the service level objective (SLO), and
\emph{minimizing instances provisioned} to each model to improve hardware utilization.
Achieving both is challenging
 due to the unpredictable short-term fluctuations in a model's instance demands,
 (5\,$\times$ required within 2 seconds),
 because the request arrival rate bursts at seconds-level~\cite{alpaserve,DBLP:conf/nsdi/0025TKS23},
 where the memory usage of each request is also unpredictable due to the
 auto-regressive nature of LLMs~\cite{alpaserve,DBLP:conf/nsdi/0025TKS23,DBLP:journals/corr/abs-2401-14351} (see also {\fig{fig:problem-statement} and \textsection{\ref{sec:problem-statement}}}).
\iffalse
For example, when serving a \TODO{xx} model under an Azure real-world serving trace~\cite{}, 
 the number of GPUs needed to process requests within SLO 
 and store intermediate results (KVCache) varies \TODO{xx} with no clear pattern ({\fig{fig:problem-statement}}). 
 These fluctuations can occur 
 due to changes in either external request rates or internal processing needs. 
 For instance, the processing time of a large language model depends on 
 how it generates results (tokens)---which is difficult to predict due to the model's auto-regressive nature. 
 \fi

 Model autoscaling is a promising solution~\cite{DBLP:conf/asplos/YangZLZLZCL22,serverless-ray,DBLP:journals/corr/abs-2401-14351,kserve}.
 With autoscaling,
  a {\maas} system only provisions the average number of instances required
  over the long term for each served model, which remains relatively stable~\cite{wang2024burstgpt}.
  This improves utilization.
  Upon bursts,
  the system automatically scales new instances,
  avoiding SLO violations due to request queueing caused by insufficient instances provisioned.

Autoscaling speed is critical in minimizing SLO violations
because the queued requests are not served until the instances are scaled.
For instance, the inference time of a Llama3-8B is 80-900\,ms on commodity GPU (A800),
while users expect a tight response time (< 1 second) for scenarios like chatbot~\cite{latency-study,DBLP:conf/osdi/ZhongLCHZL0024,aws-latency}.
Meeting such tight SLO requires less than 500\,ms scaling time,
but achieving this is challenging especially for LLMs
with 10-400\,GB parameters.
The key reason is
the \textbf{slow data plane} of autoscaling---the process of loading the model parameters to instances' GPUs.
While high bandwidth SSDs
are utilized in current work~\cite{DBLP:journals/corr/abs-2401-14351},
the speed provided by SSDs of GPU servers (2-10\,Gbps per GPU~\cite{googleinstance,awsinstance,volcanoinstance})
is still far from ideal.
For instance, loading Llama3-8\,B to a GPU takes 12.8 seconds with 10\,Gbps SSD.
Another factor obstructing fast scaling is that existing scaling methods are
\textbf{stop-the-world}:
the scaled instances cannot serve requests until all parameters are loaded.
This implies that autoscaling is directly bottlenecked by the data plane.

To mitigate the above issues, state-of-the-art systems like ServerlessLLM
further adopt a multi-tiered caching system
by caching model parameters in the host (CPU) DRAM to accelerate the data plane~\cite{DBLP:conf/osdi/GujaratiKAHKVM20,DBLP:conf/eurosys/JeongBA23}.
Under cache hit, they can leverage the fast CPU-GPU link (e.g., 256\,Gbps PCIe) to load parameters.
However, achieving a high hit rate is unfeasible:
ServerlessLLM reports a hit rate of 40--75\,\%,
which is confirmed by us (see \textsection{\ref{sec:bg-motiv}}).
The root cause is that a {\maas} typically hosts many models,
thus achieving a 100\,\% cache hit requires caching all these models on the DRAM of each host,
clearly impractical.
Vendors typically host many models because
there are hundreds of popular open-source model families designed for different purposes~\cite{huggingface}.
Meanwhile,
each model family has different scales for balancing the serving cost and accuracy~\cite{qwens}.
Finally,
developers can upload their customized fine-tuned models based on open-source models~\cite{aws-serverless-llm}.

\iffalse
so caching all models on all hosts is impractical especially when 
a model must scale across multiple instances on multiple hosts.
The hosted models come from a model zoo like
HuggingFace~\cite{huggingface}, which offers 1,170,000 available models, 
with 1,596 of them having over 15,000 downloads each.
Storing these frequently downloaded models would demand 23\,TBs of memory,
clearly exceeding the capacity of host memory and even SSDs. 
\fi

\iffalse
The state-of-the-art system---ServerlessLLM~\cite{DBLP:journals/corr/abs-2401-14351} tackles 
this problem by building a multi-tiered caching system.
It always store all models on the SSD at each host to handle host cache misses.
Unfortunately, 
we found that SSD bandwidth of the serving infrastructure is still far 
from ideal (2--10\,Gbps per GPU reported by the vendors~\cite{googleinstance,awsinstance,volcanoinstance}). 
As a result, 
we still observe a \TODO{xxx} SLO violations when serving a \TODO{xx} workload with \TODO{model}. 
\fi



\iffalse
The auto scaling time is critical to minimize SLO violations:
the state-of-the-art system ServerlessLLM~\cite{DBLP:journals/corr/abs-2401-14351}
adopts a bandwidth-optimized SSD parameter loading technique, 
but it still incurs \TODO{xx} SLO violations 
when serving a \TODO{Xx} model on the \TODO{xx} workload (\textsection{\ref{sec:bg-motiv}}). 
The scale is bottlenecked by the \textbf{data plane} of loading huge parameters to an instance: 
First, the load is \emph{lengthy}: SSD takes \TODO{xxx} seconds for loading parameters 
under common setups for serving clusters (2--10\,Gbps per GPU provided from vendors~\cite{googleinstance,awsinstance,volcanoinstance}). 
For comparison, 
the average time to first token is \TODO{xx} for \TODO{llama-2}. 
More importantly, the scale is \emph{stop}: 
only after all data have been loaded to an instance, a model can serve requests.
The stop, in combined with the lengthy load, causing SLO violations to the the queued surge requests.

A common solution is to cache the model parameters on hosts~\cite{DBLP:conf/osdi/GujaratiKAHKVM20,DBLP:conf/eurosys/JeongBA23},
which leverages the faster host-GPU link (e.g., 256\,Gbps PCIe 4.0) to speed up the data plane.
However,caching is not always feasible: 
First, a model may require scaling across multiple instances (\TODO{Xx measured}) on \TODO{xx} machines, 
but it is impractical to cache the model on all these machines.
This is because modern {\maas} systems simultaneously host many models:
Our private conversation with a top-3 cloud vendor revealed 
that thousands of models are hosted on their {\maas}. 
HuggingFace~\cite{huggingface}---a popular model zoo—contains over 1,170,000 available models, 
with 1,596 models---each with at least over 15,000 downloads---that can potentially be served by {\maas}. 
Storing these frequently used models would require 23\,TBs of memory, which clearly exceeds the host memory. 
As a result, ServerlessLLM~\cite{DBLP:journals/corr/abs-2401-14351} reports that serving suffers up to 60\% cache miss,
which is also confirmed by us. 
Second, for some serving workloads, 
even the host-GPU speed may be insufficient for efficient scaling (\textsection{\ref{sec:bg-motiv}}).
\iffalse
Thus, existing systems only adopt a conservative time-to-live caching policy~\cite{DBLP:journals/corr/abs-2401-14351},
and we found near-zero hit rate when serving the BurstGPT workload even without colocated models,
because the interval between bursts are high (\TODO{xx}). 
\fi
\fi

To achieve fast model scaling without relying on cache hit,
we make the following two key contributions:

\stitle{1. Data plane made fast with $O(1)$ or no caching with compute network multicast.\,}
First, a {\maas} is backed by fast GPU-GPU/CPU compute fabrics~\cite{nvidia-compute-network},
which are 100--400\,Gbps RDMA and even 16\,Tbps NVLink~\cite{googleinstance,awsinstance,volcanoinstance}---much
faster than SSD and comparable (or even faster than) CPU-GPU PCIe.
The compute fabric is used for data transfer during serving
and we found it largely under-utilized,
i.e., up to 7.4\,\% of total bandwidth
even in network-heavy workloads like serving LLMs with prefill and decode disaggregation~\cite{DBLP:conf/isca/PatelCZSGMB24, DBLP:conf/osdi/ZhongLCHZL0024,DBLP:journals/corr/abs-2401-11181,DBLP:conf/ppopp/CaiLMMMNS21,DBLP:journals/corr/abs-2404-09526} (\textsection{\ref{sec:bg-motiv}}).
Thus, we can borrow such fast links for accelerating the data plane of autoscaling.

Second, network-based data plane requires no or minimal caching to achieve fast scaling. 
Specifically,
if a model is already deployed on some instances,
we can directly multicast the parameters from deployed instances through the network, 
eliminating the need for caching.
Such multicast is extremely efficient
because a serial forwarding multicast~\cite{DBLP:conf/icpp/VerstoepLB96}
can load bulk data (e.g., model parameters), regardless of the number of receivers. 
Even if no instance is deployed, multicast can be done with $O(1)$ host caching by
simply broadcasting the parameters from the host with the cached model. 
This $O(1)$ caching per-model allows us to avoid all cache misses
since the aggregated host memory of all machines is sufficient to cache all models served by a {\maas}. 

\iffalse
Thus, we first propose a global parameter pool to manages 
the model parameters across all the GPU servers (\textsection{\ref{sec:design-plan}}),
and leverage network-based multicast % extensively studied~\cite{DBLP:journals/jpdc/BhatRP03,DBLP:conf/asplos/CowanMMSX23,DBLP:conf/sc/GhazimirsaeedZR20,DBLP:conf/icpp/BanikazemiMP98,DBLP:conf/dsn/BehrensJBT18} 
to accelerate the data plane of data loading. 
\fi

\iffalse
Therefore, the first challenge we need to tackle 
is to carefully coordinate different network flows to prevent interference. 
{\sys} addresses this problem by co-designing parameter loading with model dataflow 
as well as the bi-directional feature of modern networking (\textsection{\ref{sec:design-plan}}). 
\fi


Although fast networking can significantly accelerate the data plane
with minimal caching,
a stop-the-world loading remains a bottleneck in cases
when the networking is not fast enough.
For example,
to achieve at most 40\,\% SLO violations when serving a BurstGPT workload with Qwen2.5-72B,
the system needs to achieve a tight 500\,ms stop time.
Achieving so requires 576\,Gbps per-GPU\footnote{\footnotesize{72\,B model requires at least four GPUs per-instance for serving.}}
parameter transfer bandwidth,
far exceeding
the available bandwidth of typical compute network setups (e.g., 200\,Gbps per-GPU)
and even when caching at the host (256\,Gbps PCIe).
Thus, we argue that an ideal parameter loading should be \emph{live}:
before the data loading finishes, the scaled instance should be able to serve requests.

\stitle{2. Data plane made live with fine-grained scaling abstraction and cooperative execution. \,}
Model scaling cannot be live using traditional
on-demand data loading techniques commonly found in serverless computing~\cite{mitosis,DBLP:conf/eurosys/WangHW19,DBLP:conf/sosp/HuangZMLLCJLSZL24}
or inference loading overlap in PipeSwitch~\cite{pipeswitch} (\textsection{\ref{sec:overview}}),
because an instance can only emit results once all the parameters are loaded.
This stop is rooted in the coarse-grained scaling abstraction of existing systems:
they can only scale and serve at the instance level.
To realize live scaling,
our key insight is that
\emph{models can be served in a fine-grained, layer-by-layer manner}.
With this fine-grained layer-wise scaling,
we can offload part of the layer's computation from overloaded instances to scaled instances
with cooperative execution,
thus improving the overall serving throughput
even before the scaled instance has loaded all the parameters.


\iffalse
One problem when we derive the above model-aware 
cooperative execution is that it requires careful request scheduling
to decide which instance execute which layers: 
unlike a distributed serving without autoscale, 
where the load can be evenly distributed, 
the load can be imbalanced due to partially loaded data. 
As we will show, a simple best effort scaling that execute as much layers as possible
on the new instances still result in \TODO{xx} SLO violations.
Our model and load status-aware zigzag scheduling addresses this issue (\textsection{\ref{sec:design-pp}}).
\fi

\stitle{Challenges and solutions.\,}
First,
utilizing network-based multicast is non-trivial in our setup.
Though the mechanism of multicast is simple, i.e.,
simply forwarding parameters between instances with the network,
the challenges lie in generating the multicast plan,
i.e., determining how the data flows between instances.
First, we need to quickly generate an efficient plan online on
diverse network topologies since our sources and destinations are dynamically determined,
but generating an optimal plan is NP-hard on
heterogeneous networks in serving clusters~\cite{DBLP:journals/jpdc/BhatRP03}.
Second, we need to avoid network interference between the scaling and serving,
otherwise, we observed a 1.5\,$\times$ longer scale time and 50\% degraded tail TBT (\textsection{\ref{sec:overview}}).
Current solutions~\cite{DBLP:conf/asplos/CowanMMSX23,DBLP:conf/ppopp/CaiLMMMNS21,DBLP:conf/sc/GhazimirsaeedZR20} mainly
target offline scenarios like training,
so they can tolerate long plan generation time and
don't need to consider interference from serving workloads.
To address the issue, we propose a
model-aware multicast planner,
which leverages the key features of compute network and the
static data flow in model serving to quickly generate a near-optimal,
interference-free multicast plan for scaling (\textsection{\ref{sec:design-plan}}).

Second, it is challenging to schedule how requests are executed
between deployed and live scaling instances,
i.e., which instance executes which layers.
The challenge lies in the fact that the serving capability of the scaling instances
is limited---it can only execute layers with parameters loaded,
and this capability is dynamically changing.
A naive best-effort scaling that executes as many layers as possible
cannot balance the load because at the beginning of autoscaling,
the new instances can only execute few layers, 
with requests still queued at the overloaded instances.
A better solution is to adjust the load holistically by considering future incoming layers,
and we realize this with a ZigZag pipeline scheduling
and achieve 50\% tail latency reduction under bursty workloads
(\textsection{\ref{sec:design-pp}}).


\stitle{Demonstration with {\sys}. \,}
We built {\sys}, an {\maas} system with the fastest autoscaling speed with $O(1)$ caching.
We adopted a global parameter pool to cache the model parameters across all the machines,
and integrated the aforementioned interference-free multicast plan for scaling
and efficient ZigZag scheduling-based live scheduling.
To show the effectiveness of {\sys}, 
we evaluated {\sys} by running real-world traces
(i.e., BurstGPT~\cite{wang2024burstgpt}, AzureCode and AzureConv~\cite{AzurePublicDataset})
across a variety of recent models
with different sizes and architectures, including Llama3-8B, Mistral-24B, and Qwen2.5-72B.
First, {\sys} has 47--75\,\% shorter time-to-first token,
and has up to 94\,\% shorter time-between-tokens 
than the state-of-the-art work (ServerlessLLM~\cite{DBLP:journals/corr/abs-2401-14351}). 
Second, compared to serving systems without autoscaling support,
i.e., vLLM~\cite{vllm} and DistServe~\cite{DBLP:conf/osdi/ZhongLCHZL0024},
{\sys} reduces the GPU used for serving a single model by 49\,\%
with no SLO violations compared to an over-provisioning setup
that provisions the GPUs based on the maximum request rate.
{\sys} is open-sourced at \burl{https://github.com/blitz-serving/blitz-scale}. 
